% Detailed derivation gathered at the end of the owning structure.

\begin{derivation}[从\eqnrefs{eq:free-generating-functional-r3}到\eqnrefs{eq:qft-QJ-second-derivative-dp6}]

正文\eqnrefs{eq:free-generating-functional-r3} 把自由理论的全部源依赖写成一个二次指数。二次高斯源指数的两次源微分恰好产生一个 Feynman 传播子。这里计算这一二阶导数；四点及更高点则由同一高斯指数的高阶微分产生，并由 Wick 配对统一组织。

记
\begin{equation*}
 \mathscr Q[\mathcal J]
 \equiv
 -\frac{1}{2\hbar^2c^2}
 \int\dd^4 x\,\dd^4 y\,
 \mathcal J(x)\Delta_F(x-y)\mathcal J(y),
 \qquad
 \mathcal Z_0[\mathcal J]=\ee^{\mathscr Q[\mathcal J]}.
\end{equation*}
由于实标量场的 Feynman 核满足
$\Delta_F(x-y)=\Delta_F(y-x)$，第一次泛函导数为
\begin{align*}
 \frac{\delta\mathscr Q}{\delta\mathcal J(x_1)}
 &=-\frac{1}{2\hbar^2c^2}
 \int\dd^4 y\,\Delta_F(x_1-y)\mathcal J(y)
 -\frac{1}{2\hbar^2c^2}
 \int\dd^4 x\,\mathcal J(x)\Delta_F(x-x_1)\\
 &=-\frac{1}{\hbar^2c^2}
 \int\dd^4 y\,\Delta_F(x_1-y)\mathcal J(y).
\end{align*}
于是
\begin{equation*}
 \frac{\delta\mathcal Z_0}{\delta\mathcal J(x_1)}
 =\mathcal Z_0
 \frac{\delta\mathscr Q}{\delta\mathcal J(x_1)}.
\end{equation*}
再对 $\mathcal J(x_2)$ 求导，乘积法则给出
\begin{align*}
 \frac{\delta^2\mathcal Z_0}
 {\delta\mathcal J(x_1)\delta\mathcal J(x_2)}
 =&\;\mathcal Z_0
 \frac{\delta\mathscr Q}{\delta\mathcal J(x_1)}
 \frac{\delta\mathscr Q}{\delta\mathcal J(x_2)}
 +\mathcal Z_0
 \frac{\delta^2\mathscr Q}
 {\delta\mathcal J(x_1)\delta\mathcal J(x_2)},\\
 \frac{\delta^2\mathscr Q}
 {\delta\mathcal J(x_1)\delta\mathcal J(x_2)}
 =&-\frac{1}{\hbar^2c^2}\Delta_F(x_1-x_2),
\end{align*}
这就是正文\eqnrefs{eq:qft-QJ-second-derivative-dp6}。在零源处，两个一阶导数都含显式 $\mathcal J$ 而消失，所以剩下的二阶核直接给出自由二点函数；更高点的配对结构同样来自这一高斯性质，并可整理成 Wick 配对。
\end{derivation}
